\documentclass[11pt, UTF8]{ctexart}
\usepackage{amsmath, amsfonts, amssymb}
\usepackage{extarrows} 
\usepackage{geometry} 
\usepackage{xcolor} 
\usepackage{enumitem} 
\usepackage[most]{tcolorbox} 
\usepackage{fancyhdr} 
\usepackage{diagbox} 
\usepackage{tikz}    
\setlength{\headheight}{13.6pt}
\geometry{a4paper, margin=0.8in} 
\pagestyle{fancy}
\fancyhf{}
\lhead{\color{fblue} \hwdate 作业} 
\rhead{\color{fblue} \today}
\cfoot{\thepage}

\definecolor{fblue}{RGB}{0, 74, 142}
\definecolor{fred}{RGB}{204, 26, 26}

\newtcolorbox{problem}[2]{
    enhanced,
    colback=white,
    colframe=fblue!80,
    arc=2mm,
    boxrule=0.8pt,
    before upper={\textbf{\color{fblue} 问题 #1 (#2)}\hspace{0.5em}},
    before skip=1.5em,
    after skip=1.5em,
    left=3mm, right=3mm, top=2mm, bottom=2mm
}

\newcommand{\hwdate}{6月4日} 
\newcommand{\Prob}{P}
\newcommand{\comp}[1]{\overline{#1}\mskip 3mu}
\newcommand{\ans}[1]{{\color{fred}\boldsymbol{#1}}}

\begin{document}

\begin{center}
    {\LARGE \bfseries \color{fblue} 作业}
\end{center}

% ---------------- 题目 1 ----------------
\begin{problem}{1}{7. 4}
\textbf{(1)} 设总体 \(X\) 具有分布律
\begin{center}
\begin{tabular}{c|ccc}
\(X\) & \(1\) & \(2\) & \(3\) \\ \hline
\(p_k\) & \(\theta^2\) & \(2\theta(1-\theta)\) & \((1-\theta)^2\)
\end{tabular}
\end{center}
其中 \(\theta(0<\theta<1)\) 为未知参数. 已知取得了样本值 \(x_1=1, x_2=2, x_3=1\). 试求 \(\theta\) 的矩估计值和最大似然估计值.

\textbf{(2)} 设 \(X_1, X_2, \cdots, X_n\) 是来自参数为 \(\lambda\) 的泊松分布总体的一个样本, 试求 \(\lambda\) 的最大似然估计量及矩估计量.

\textbf{(3)} 设随机变量 \(X\) 服从以 \(r, p\) 为参数的负二项分布, 其分布律为
\[ P\{X=x_k\} = \binom{x_k-1}{r-1} p^r (1-p)^{x_k-r}, \quad x_k = r, r+1, \cdots, \]
其中 \(r\) 已知, \(p\) 未知. 设有样本值 \(x_1, x_2, \cdots, x_n\), 试求 \(p\) 的最大似然估计值.
\end{problem}

\paragraph{\color{fblue}解}
\textbf{(1) (i)} 矩估计法：
\[ \mu_1 = \theta^2 + 2 \cdot 2\theta(1-\theta) + 3 \cdot (1-\theta)^2 = 3 - 2\theta. \]
解得 \(\theta = \frac{1}{2}(3 - \mu_1)\), 故得 \(\theta\) 的矩估计值为
\[ \hat{\theta} = \frac{1}{2}(3 - \bar{x}). \]
令 \(\bar{x} = \frac{1}{3}(x_1 + x_2 + x_3) = \frac{1}{3}(1+2+1) = \frac{4}{3}\), 故 \(\theta\) 的矩估计值为 \(\ans{\hat{\theta} = \frac{5}{6}}\).

\textbf{(ii)} 最大似然估计法：由给定的样本值, 得似然函数为
\begin{align*}
L &= \prod_{i=1}^3 P\{X_i = x_i\} = P\{X_1 = 1\}P\{X_2 = 2\}P\{X_3 = 1\} \\
  &= \theta^2 \cdot 2\theta(1-\theta) \cdot \theta^2 = 2\theta^5(1-\theta),
\end{align*}
\[ \ln L = \ln 2 + 5\ln \theta + \ln(1-\theta). \]
令
\[ \frac{d}{d\theta}\ln L = \frac{5}{\theta} - \frac{1}{1-\theta} = 0, \]
得 \(\theta\) 的最大似然估计值为 \(\ans{\hat{\theta} = \frac{5}{6}}\).

\vspace{1em}
\textbf{(2)} (i) 设 \(x_1, x_2, \cdots, x_n\) 是相应于样本 \(X_1, X_2, \cdots, X_n\) 的样本值, 则似然函数为
\[ L = \prod_{i=1}^n \frac{\lambda^{x_i}}{x_i!} e^{-\lambda} = e^{-n\lambda} \lambda^{\sum_{i=1}^n x_i} \bigg/ \prod_{i=1}^n x_i!, \]
\[ \ln L = -n\lambda + \left(\sum_{i=1}^n x_i\right) \ln \lambda - \ln \prod_{i=1}^n x_i!. \]
令
\[ \frac{d}{d\lambda}\ln L = -n + \left(\sum_{i=1}^n x_i\right) \bigg/ \lambda = 0, \]
得 \(\lambda\) 的最大似然估计值为 \(\hat{\lambda} = \frac{1}{n}\sum_{i=1}^n x_i = \bar{x}\), 最大似然估计量为 \(\ans{\hat{\lambda} = \bar{X}}\).

(ii) 因 \(\mu_1 = E(X) = \lambda\), 故 \(\lambda\) 的矩估计量也是 \(\ans{\hat{\lambda} = \bar{X}}\).

\vspace{1em}
\textbf{(3)} 似然函数为
\begin{align*}
L(p) &= \prod_{k=1}^n P\{X = x_k\} = \prod_{k=1}^n \binom{x_k-1}{r-1} p^r (1-p)^{x_k-r} \\
     &= \left[\prod_{k=1}^n \binom{x_k-1}{r-1}\right] p^{nr} (1-p)^{\sum_{k=1}^n x_k - nr},
\end{align*}
\[ \ln L = \ln C + nr \ln p + \left(\sum_{k=1}^n x_k - nr\right) \ln(1-p), \quad C \text{为常数}. \]
令
\[ \frac{d}{dp}(\ln L) = \frac{nr}{p} - \frac{1}{1-p}\left(\sum_{k=1}^n x_k - nr\right) = 0, \]
得 \(p\) 的最大似然估计值为
\[ \ans{\hat{p} = \frac{r}{\bar{x}}}. \]

\vspace{1.5em}

% ---------------- 题目 2 ----------------
\begin{problem}{2}{7. 5}
设某种电子器件的寿命(以 h 计) \(T\) 服从双参数的指数分布, 其概率密度为
\[ f(t) = \begin{cases}
\frac{1}{\theta} \mathrm{e}^{-(t-c)/\theta}, & t \ge c, \\
0, & \text{其他}.
\end{cases} \]
其中 \(c, \theta \, (c, \theta > 0)\) 为未知参数. 自一批这种器件中随机地取 \(n\) 件进行寿命试验. 设它们的失效时间依次为 \(x_1 \le x_2 \le \cdots \le x_n\).
\begin{enumerate}[label=(\arabic*)]
    \item 求 \(c\) 与 \(\theta\) 的最大似然估计值.
    \item 求 \(c\) 与 \(\theta\) 的矩估计量.
\end{enumerate}
\end{problem}

\paragraph{\color{fblue}解}
\textbf{(1)} 似然函数为
\[ L(\theta, c) = L(x_1, x_2, \cdots, x_n; \theta, c)
= \begin{cases}
\prod_{i=1}^n \frac{1}{\theta} \mathrm{e}^{-(x_i-c)/\theta}, & x_1, x_2, \cdots, x_n \ge c, \\
0, & \text{其他}.
\end{cases} \]
由题设 \(x_1 \le x_2 \le \cdots \le x_n\), 故 \(x_1, x_2, \cdots, x_n \ge c\) 相当于 \(x_1 \ge c\), 因而上式相当于
\[ L(\theta, c) = \begin{cases}
\frac{1}{\theta^n} \mathrm{e}^{-\left(\sum_{i=1}^n x_i - nc\right)/\theta}, & c \le x_1, \\
0, & c > x_1.
\end{cases} \]
可知当 \(c \le x_1\) 时, \(L(\theta, c)\) 随 \(c\) 的增加而递增, 而当 \(c > x_1\) 时 \(L(\theta, c) = 0\), 因而对于固定的 \(\theta, L(\theta, c)\) 在 \(c = x_1\) 取到最大值, 从而知应取 \(\hat{c} = x_1\).

另外, 当 \(c \le x_1\) 时, 将 \(L(\theta, c)\) 取自然对数得
\[ \ln L(\theta, c) = -n\ln \theta - \frac{1}{\theta}\left(\sum_{i=1}^n x_i - nc\right). \]
令 \(\frac{\partial}{\partial \theta}\ln L(\theta, c) = 0\), 得
\[ \frac{\partial}{\partial \theta}\ln L(\theta, c) = -\frac{n}{\theta} + \frac{\sum_{i=1}^n x_i - nc}{\theta^2} = 0, \]
于是
\[ \theta = \bar{x} - c. \]
由此可知 \(c, \theta\) 的最大似然估计值为
\[ \ans{\begin{cases}
\hat{c} = x_1, \\
\hat{\theta} = \bar{x} - x_1.
\end{cases}} \]

\vspace{1em}
\textbf{(2)} \(\mu_1 = \int_{-\infty}^{+\infty} t f(t) dt = \int_c^{+\infty} \frac{t}{\theta} \mathrm{e}^{-(t-c)/\theta} dt,\)
令 \(u = \frac{t-c}{\theta}\), 得
\[ \mu_1 = \int_0^{+\infty} (\theta u + c) \mathrm{e}^{-u} du = c + \theta \Gamma(2) = c + \theta. \]
\[ \mu_2 = \int_{-\infty}^{+\infty} t^2 f(t) dt = \int_c^{+\infty} \frac{t^2}{\theta} \mathrm{e}^{-(t-c)/\theta} dt, \]
令 \(u = \frac{t-c}{\theta}\), 得
\[ \mu_2 = \int_0^{+\infty} (\theta u + c)^2 \mathrm{e}^{-u} du = \theta^2 \Gamma(3) + 2c\theta \Gamma(2) + c^2 = 2\theta^2 + 2c\theta + c^2 = (c+\theta)^2 + \theta^2, \]
由此得
\[ \theta = \sqrt{\mu_2 - \mu_1^2}, \quad c = \mu_1 - \sqrt{\mu_2 - \mu_1^2}. \]
将上两式中的 \(\mu_1, \mu_2\) 分别换成 \(A_1 = \bar{X}, A_2 = \frac{1}{n}\sum_{i=1}^n X_i^2\), 并注意到 \(A_2 - A_1^2 = \frac{1}{n}\sum_{i=1}^n (X_i - \bar{X})^2\), 就得到 \(\theta\) 及 \(c\) 的矩估计量为
\[ \ans{\hat{\theta} = \sqrt{\frac{1}{n}\sum_{i=1}^n (X_i - \bar{X})^2}, \quad \hat{c} = \bar{X} - \sqrt{\frac{1}{n}\sum_{i=1}^n (X_i - \bar{X})^2}}. \]

\vspace{1.5em}

% ---------------- 题目 3 ----------------
\begin{problem}{3}{7. 15}
设有 \(k\) 台仪器, 已知用第 \(i\) 台仪器测量时, 测定值总体的标准差为 \(\sigma_i \, (i=1, 2, \cdots, k)\). 用这些仪器独立地对某一物理量 \(\theta\) 各观察一次, 分别得到 \(X_1, X_2, \cdots, X_k\). 设仪器都没有系统误差, 即 \(E(X_i) = \theta \, (i=1, 2, \cdots, k)\). 问 \(a_1, a_2, \cdots, a_k\) 取何值, 方能使使用 \(\hat{\theta} = \sum_{i=1}^k a_i X_i\) 估计 \(\theta\) 时, \(\hat{\theta}\) 是无偏的, 并且 \(D(\hat{\theta})\) 最小?
\end{problem}

\paragraph{\color{fblue}解}
要使 \(\hat{\theta}\) 是 \(\theta\) 的无偏估计, 则必须
\[ \theta = E(\hat{\theta}) = E(a_1 X_1 + a_2 X_2 + \cdots + a_k X_k) = (a_1 + a_2 + \cdots + a_k)\theta. \]
即必须
\[ a_1 + a_2 + \cdots + a_k = 1. \quad (*_1) \]
又由题设知 \(D(X_i) = \sigma_i^2, i=1, 2, \cdots, k\), 且 \(X_1, X_2, \cdots, X_k\) 相互独立, 故有
\begin{align*}
D(\hat{\theta}) &= D\left(\sum_{i=1}^k a_i X_i\right) = \sum_{i=1}^k D(a_i X_i) = \sum_{i=1}^k a_i^2 D(X_i) \\
&= a_1^2 \sigma_1^2 + a_2^2 \sigma_2^2 + \cdots + a_k^2 \sigma_k^2.
\end{align*}
为求 \(D(\hat{\theta})\) 在条件 \((*_1)\) 下的最小值, 用拉格朗日乘数法, 作函数
\[ g(a_1, a_2, \cdots, a_k, \lambda) = a_1^2 \sigma_1^2 + a_2^2 \sigma_2^2 + \cdots + a_k^2 \sigma_k^2 + \lambda(a_1 + a_2 + \cdots + a_k - 1). \]
令
\[ \frac{\partial g}{\partial a_1} = 2a_1 \sigma_1^2 + \lambda = 0, \quad \frac{\partial g}{\partial a_2} = 2a_2 \sigma_2^2 + \lambda = 0, \cdots, \]
\[ \frac{\partial g}{\partial a_k} = 2a_k \sigma_k^2 + \lambda = 0, \quad \frac{\partial g}{\partial \lambda} = \sum_{i=1}^k a_i - 1 = 0. \]
得
\[ a_1 = -\frac{\lambda}{2\sigma_1^2}, \quad a_2 = -\frac{\lambda}{2\sigma_2^2}, \cdots, \quad a_k = -\frac{\lambda}{2\sigma_k^2}, \quad \sum_{i=1}^k a_i = 1. \quad (*_2) \]
将前 \(k\) 个式子代入末式, 得
\[ -\frac{\lambda}{2}\sum_{i=1}^k \frac{1}{\sigma_i^2} = 1. \]
在上式中记 \(\sum_{i=1}^k \frac{1}{\sigma_i^2} = \frac{1}{\sigma_0^2}\), 即 \(-\frac{\lambda}{2} = \sigma_0^2\), 代入 \((*_2)\) 式, 得
\[ a_1 = \frac{\sigma_0^2}{\sigma_1^2}, \quad a_2 = \frac{\sigma_0^2}{\sigma_2^2}, \cdots, \quad a_k = \frac{\sigma_0^2}{\sigma_k^2}. \]
即当 \(\ans{a_i = \frac{\sigma_0^2}{\sigma_i^2}}, i=1, 2, \cdots, k\), 用 \(\hat{\theta} = \sum_{i=1}^k a_i X_i\) 估计 \(\theta\) 时, \(\hat{\theta}\) 是无偏的且其方差最小(相对于其他无偏估计来说).

\end{document}